% !PRINTED 234
\weberpage{234}
%<<CONT

\gesperrt{3. Wenn} $J(x_1, x_2)$ \gesperrt{irgend eine invariante Form der Ikosa\"edergruppe ist, und} $\xi$ \gesperrt{eine ihrer Wurzeln, so dass} $J(\xi, 1) = 0$ \gesperrt{ist, so sind auch alle die Gr\"ossen Wurzeln von} $J$, \gesperrt{die aus} $\xi$ \gesperrt{durch die gebrochenen Ikosa\"edersubstitutionen hervorgehen.}

Wenn daher $J$ mit einer der Functionen $f$, $H$, $T$ einen Theiler gemein hat, so ist $J$ durch die betreffende Form theilbar.

Wir denken uns nun zun\"achst aus $J$ m\"oglichst hohe Potenzen der drei Grundformen $f$, $H$, $T$ weggehoben, so dass $J$ zu diesen Functionen theilerfremd ist. Ist dann $\xi$ eine Wurzel von $J$, so k\"onnen wir $\eta$ so bestimmen, dass $\xi$ eine Wurzel der Form
\[
\Theta = H^3 - \eta f^5
\]
ist, und $\xi$ ist dann nach 2. eine einfache Wurzel von $\Theta$. Es m\"ussen dann auch nach 3. alle \"ubrigen Wurzeln von $\Theta$ zugleich Wurzeln von $J$ sein, d. h. $J$ ist durch $\Theta$ theilbar. Der Quotient ist wieder eine invariante Form des Ikosa\"eders, und der Schluss l\"asst sich wiederholen. Wir kommen so zu dem Satze:

\gesperrt{4. Jede Invariante des Ikosa\"eders l\"asst sich in der Form darstellen:}
\weq{2}{J(x_1,x_2) = C f^\alpha H^\beta T^\gamma F(H^3, f^5),}
\gesperrt{worin $\alpha$, $\beta$, $\gamma$ nicht negative ganze Zahlen, $C$ eine Constante und $F$ eine ganze homogene Function bedeuten.}

\begin{center}
\S.\,62.\\
\gesperrt{Poly\"edergruppen der zweiten Art. Krystallographische Gruppen.}
\end{center}

Wir haben schon im \S.\,51 auf Gruppen linearer tern\"arer Substitutionen aufmerksam gemacht, die ausser den eigentlichen auch uneigentlich orthogonale Substitutionen enthalten, und die wir \gesperrt{Gruppen der zweiten Art} nennen.

Wir wollen die endlichen unter ihnen jetzt als \gesperrt{Poly\"edergruppen} der zweiten Art oder auch als \gesperrt{erweiterte Poly\"edergruppen}, und die darin enthaltenen Substitutionen mit der
%CONT>>
